% No preamble. Unique prefix ob-. Follows the oi- overlap theorem.
% Strict background support under the explicit local-precision hypotheses.

\section{Strict background support for overlapping interactions}
\label{ob-section}

The preceding expansion can also be made strictly local in its
background variables. This requires a background-independent baseline
and a small local precision decomposition, in addition to local
background dependence of the interactions. The next statement checks
these additional inputs explicitly.

\subsection{Baseline, local precision and constants}

Let $\Lambda$ be a finite subset of $\mathbb Z^d$ with the ambient
lattice graph metric and fixed cell dimension $D$. The background
domain is a product of complex coordinate domains. Write
\begin{equation}\label{ob-precision}
 P(A)=P_0+\sum_Y B_Y(A|_Y),\qquad C_0=P_0^{-1},
\end{equation}
where $C_0$ is complex symmetric and independent of every background variable, each
$B_Y$ is complex symmetric, has rows and columns supported in $Y$,
and reads only $A|_Y$. Assume
\begin{equation}\label{ob-spectrum}
 \operatorname{Herm}C_0\ge\lambda I,\qquad
 \|C_0\|\le\Lambda_0,\qquad 0<\lambda\le\Lambda_0,
\end{equation}
and, for fixed $0<\kappa<m$,
\begin{align}
 K&=\sup_b\sum_c e^{(\kappa+2m)d(b,c)}
                         \|(C_0)_{bc}\|<\infty,\notag\\
 \nu&=\sup_b\sum_{Y\ni b}\sum_{c\in Y}
       e^{(\kappa+2m)\ell(Y)}
                         \|(B_Y)_{bc}\|_\infty<\infty. \label{ob-norms}
\end{align}
Symmetry gives the same column bounds. All constants are uniform
in volume and background. Define
\begin{align}
 \delta&=\lambda/4,&
 V_*&=8\Lambda_0^2/\lambda,&
 J_{bc}&=K e^{-m d(b,c)}\quad(b\ne c),\notag\\
 A_\kappa&=K\sum_{z\ne0}e^{-(m-\kappa)|z|},&
 \bar K&=K+\delta,& a&=\lambda/2,\notag\\
 s&=\min\{m/2,\;aem/(4\bar K)\},&
 L&=\frac4\lambda\sum_{z\in\mathbb Z^d}e^{-s|z|}. \label{ob-baseline-constants}
\end{align}
Set $J_{bb}=0$. Require $K\nu<1$ and put
\begin{equation}\label{ob-perturbation-constants}
 v=\frac{K^2\nu}{1-K\nu},\qquad
 v_1=v(1+\delta/K),\qquad
 u=Lv_1,\qquad w=\frac{L^2v_1}{1-u}.
\end{equation}
The sufficient precision smallness conditions are
\begin{equation}\label{ob-precision-smallness}
 v_1\le\lambda/4,\qquad u\le1/2,\qquad w\le1/(2V_*).
\end{equation}
They are chosen before increasing the cutoff radius.
The covariance is $C(A)=P(A)^{-1}$, given by its convergent
Neumann series.

Retain the fixed cutoff and interaction multiplier hypotheses
of \eqref{oi-multiplier}, now requiring each $V_Y$ to read
background only in $Y$. Define
\begin{align}
 \widehat K_G&=(8V_*/\lambda)^{D/2},&
 \widehat p_R&=2^{D/2}e^{-R^2/(8V_*)},&
 \widehat\tau&=\widehat K_G\widehat p_R,\notag\\
 \widehat\epsilon&=
  \sup_i\sum_{Y\ni i}
   v_Y\widehat K_G^{|Y|}e^{4|Y|+\kappa\ell(Y)},&
 \widehat\eta&=e^3\widehat\tau+\widehat\epsilon,&
 \widehat\eta_*&=\frac1{4\max(1,e^2A_\kappa/\delta)}.
                                                        \label{ob-smallness}
\end{align}

\begin{theorem}[Strictly local overlapping-interaction activities]
\label{ob-theorem}
Under \eqref{ob-precision}--\eqref{ob-precision-smallness} and
$\widehat\eta\le\widehat\eta_*$, the logarithm has a convergent
representation by connected lattice supports:
\begin{equation}\label{ob-local-log}
 \log Z_V(A)=\sum_H E_{\Lambda,H}(A|_H),\qquad
 \sup_b\sum_{H\ni b}e^{a_*|H|}
        \|E_{\Lambda,H}\|_\infty
 \le2S_{\kappa/2}\widehat\eta,
\end{equation}
where
\[
 a_*=\min(1/2,\kappa/4),\qquad
 S_{\kappa/2}=\sum_{z\in\mathbb Z^d}e^{-(\kappa/2)|z|}.
\]
Each activity is holomorphic and reads only the indicated
background coordinates. Supports may lie in the lattice hull of
$\Lambda$. If
\begin{equation}\label{ob-quarter}
 e^3\widehat\tau\le\widehat\eta_*/4,\qquad
 \widehat\epsilon\le\widehat\eta_*/4,
\end{equation}
the same canonical choice obeys
\begin{equation}\label{ob-difference}
 \sup_b\sum_{H\ni b}e^{a_*|H|}
  \|E_{\Lambda,H,V}-E_{\Lambda,H,0}\|_\infty
 \le8S_{\kappa/2}\widehat\epsilon.
\end{equation}
\end{theorem}

\subsection{The support algebra and the actual marginal}

For a nonempty moment support $M$, let $\ell_M(X)$ be the
minimum total length of a forest on $M\cup X$ whose every
component meets $M$. The roots may remain in different components.
For formal local families $F=(F_X)_X$, define
\begin{equation}\label{ob-algebra}
 \|F\|_{\mathcal B_M}
   =\sum_Xe^{\kappa\ell_M(X)}\|F_X\|_\infty,\qquad
 (FG)_Z=\sum_{X\cup X'=Z}F_XG_{X'}.
\end{equation}
Every $F_X$ is holomorphic in $A|_X$. The product-domain
assumption gives its uniform local-function norm.
The inequality
$\ell_M(X\cup X')\le\ell_M(X)+\ell_M(X')$
makes this a Banach algebra. Enlarging the root set is
contractive. Evaluation, obtained by summing the local
coefficients, is a bounded algebra homomorphism.

Expand $C=C_0-C_0BC_0+C_0BC_0BC_0-\cdots$ in the individual
precision labels $B_Y$. Retain these labels before grouping by
the union of their supports. For endpoints $b,c\in M$, every
term connects all its background labels to its endpoints by
covariance links and trees inside the supports $Y$. If their
total length is $L_{\rm path}$, both
$\ell_M(\cup Y)\le L_{\rm path}$ and
$d(b,c)\le L_{\rm path}$. Successive weighted matrix sums give
$K^{j+1}\nu^j$ at order $j$. Therefore, for
$\Delta C=C-C_0$,
\begin{equation}\label{ob-covariance-algebra}
 \max\left\{
  \sup_{b\in M}\sum_{c\in M}e^{2md(b,c)}
                         \|(\Delta C)_{bc}\|_{\mathcal B_M},
  \sup_{c\in M}\sum_{b\in M}e^{2md(b,c)}
                         \|(\Delta C)_{bc}\|_{\mathcal B_M}
 \right\}\le v.
\end{equation}
This bounds the whole covariance marginal, including its
background dependence, rather than only visible forest edges.

\subsection{Molecule masks and cell endpoint polydiscs}

At a molecule BKAR point, its induced cell mask $\sigma$
is one inside every primitive molecule and has the BKAR values
between molecules. Each molecule partition mask is exactly
a cell partition mask whose blocks are unions of molecules.
Thus $\sigma$ is in the convex hull of cell partition masks
and $0\le\sigma_{bc}\le1$.

Choose one cell endpoint pair for each molecule-tree edge,
with cell degrees $d_b$, and use
$|z_{bc}|<\delta/[\max(d_b,d_c)J_{bc}]$.
The endpoint graph need not span or connect every cell.
Let $M_0$ be the resulting masked and perturbed baseline.
The row-sum argument gives
\begin{equation}\label{ob-m0}
 \operatorname{Herm}M_0\ge\lambda I/2,\qquad
 \|M_0\|\le2\Lambda_0,\qquad
 \|M_0\|_{\mathrm{Schur},m}\le K+\delta.
\end{equation}
For the last bound, each selected edge satisfies
\[
 e^{md(b,c)}|z_{bc}|\|(C_0)_{bc}\|
 \le\frac{\delta}{\max(d_b,d_c)}
                     e^{-\kappa d(b,c)}.
\]
Its row sum is at most $\delta$, using only the number of
incident selected edges. No spanning-tree property on the
cell set has been used.

Conjugation by $D_{bb}=e^{s d(b,c)}$ changes $M_0$ in
operator norm by at most
$2s(K+\delta)/(em)\le a/2$. Since its Hermitian part is at
least $a$, the inverse entries are bounded by
$(2/a)e^{-sd(b,c)}$, and hence
\begin{equation}\label{ob-inverse-schur}
 \|M_0^{-1}\|_{\mathrm{Schur}}\le L.
\end{equation}

Apply the same mask and increments to $\Delta C$, obtaining
an algebra-valued matrix $E$. The mask part has norm at most
$v$. By \eqref{ob-covariance-algebra}, the selected-edge part
has row and column sums at most $\delta v/K$:
\[
 |z_{bc}|v e^{-2md(b,c)}
   \le\frac{\delta v}{K\max(d_b,d_c)}e^{-md(b,c)}.
\]
Consequently
\begin{equation}\label{ob-e}
 \|E\|_{\mathrm{Schur}(\mathcal B_M)}\le v_1.
\end{equation}
These are precisely the estimates needed for the density
expansion, and they remain valid for the molecule masks and
the original cell endpoint polydiscs.

\subsection{The entire density, including determinant normalization}

Put $F=M_0^{-1}E$. In the matrix support algebra,
\[
 (M_0+E)^{-1}-M_0^{-1}
       =\sum_{j\ge1}(-F)^jM_0^{-1},
\]
whose Schur norm is at most $w$, while
\[
 \|\operatorname{Tr}\log(I+F)\|_{\mathcal B_M}
       \le D|M|[-\log(1-u)].
\]
The exact normalized density identity is
\begin{equation}\label{ob-density-identity}
 p_{M_0+E}(x)=p_{M_0}(x)
 \exp\left[
 -\tfrac12\operatorname{Tr}\log(I+F)
 -\tfrac12x^T\big((M_0+E)^{-1}-M_0^{-1}\big)x
 \right].
\end{equation}
Taking its algebra norm, the determinant factor costs at most
$2^{D|M|/2}$. The quadratic factor costs at most
$e^{w\|x\|^2/2}$, which by
\eqref{ob-precision-smallness} uses at most half of the
baseline Gaussian decay. It follows that
\begin{equation}\label{ob-density-bound}
 \|p_{M_0+E}(x)\|_{\mathcal B_M}
       \le\widehat K_G^{|M|}p_{2V_*I}(x_M).
\end{equation}
Here $p_{2V_*I}$ is the independent product Gaussian density
with that covariance in every original cell. All inverse,
logarithm, exponential and support series are absolutely
convergent under the displayed majorants.

For an interaction supported on $Y\subset M$,
$\ell_M(Y)=0$. Thus $w_Y$ is an algebra element with
norm at most $v_Y$. A primitive interaction molecule on $S$
has norm at most $M_S$ from \eqref{oi-molecules}; a bad
singleton has empty background support.
Integrating \eqref{ob-density-bound} against their product
gives on the entire endpoint polydisc
\begin{equation}\label{ob-moment-bound}
 \left\|\mathbb E_{M_0+E}\prod_a f_{S_a}\right\|_{\mathcal B_M}
 \le\prod_a\left(
   \widehat K_G^{|S_a|}M_{S_a}
       +\widehat\tau\mathbf1_{\{|S_a|=1\}}\right).
\end{equation}

\subsection{Factorization with overlapping background labels}

Retain independent formal scalar tags on the precision labels,
and on the interaction labels if desired, until absolute
convergence is proved. At a forest mask the complete marginal
Gaussian matrix is block diagonal over the forest components.
Both its determinant and inverse therefore factorize in
\eqref{ob-density-identity}, giving an identity of formal
power series. This remains true when different components
read overlapping background coordinates. Grouping coefficients
by the union of their background labels respects multiplication.
The density for a component is its actual marginal density,
with its own root set; injection into the algebra with all
moment roots is contractive.

Thus forest factorization holds in the support representation
itself, not merely after evaluating $A$.
Banach-valued Cauchy applied to \eqref{ob-moment-bound} yields
the same cell-factorial bound \eqref{oi-tree-bound}.
The overlap-animal and planted-tree estimates give
\begin{equation}\label{ob-zeta}
 \sup_i\sum_{M\ni i,X}
 e^{\kappa\ell(M)+\kappa\ell_M(X)+2|M|}
            \|\zeta(M;X)\|_\infty\le\widehat\eta.
\end{equation}
Hard-core compatibility is still disjointness of the marked
supports $M$, and does not require disjoint background supports.

For a connected Mayer cluster, the root and union inequalities
give
\[
 \ell(\cup M_a)\le\sum_a\ell(M_a),\qquad
 \ell_{\cup M_a}(\cup X_a)
                  \le\sum_a\ell_{M_a}(X_a).
\]
The same pinned cluster estimate therefore yields
\begin{equation}\label{ob-marked-log}
 \log Z_V=\sum_{M\ne\varnothing,X}U(M;X)(A|_X),\qquad
 \sup_i\sum_{M\ni i,X}
 e^{\kappa[\ell(M)+\ell_M(X)]+|M|}
                 \|U(M;X)\|_\infty\le2\widehat\eta.
\end{equation}

\subsection{Re-rooting at the final connected background support}

Choose deterministic minimizing trees for $M$, minimizing
forests attaching $X$ to $M$, and shortest lattice paths for
their edges. Their union $H(M,X)$ is connected and contains
$M\cup X$. With $L_*=\ell(M)+\ell_M(X)$,
\[
 |H(M,X)|\le |M|+L_*,\qquad
 d(a,b)\le L_*\quad(a\in M,\ b\in H(M,X)).
\]
Consequently
\begin{equation}\label{ob-hull}
 e^{a_*|H|}
 \le e^{|M|+\kappa L_*}e^{-(\kappa/2)d(a,b)}.
\end{equation}
Group $U(M;X)$ by this fixed hull to define $E_{\Lambda,H}$.
To root a sum at $b\in H$, overcount by all $a\in M$ and
use \eqref{ob-marked-log}; the lattice convolution costs
only $S_{\kappa/2}$. This proves \eqref{ob-local-log}.

Under \eqref{ob-quarter}, the multiplier interpolation
$w_Y\mapsto tw_Y$ is holomorphic in the same support algebra.
The Banach-valued Cauchy estimate of the preceding section
gives $8\widehat\epsilon$ in the marked rooted norm.
The fixed hull map has norm at most $S_{\kappa/2}$ by
\eqref{ob-hull}, proving \eqref{ob-difference}.

\subsection{What the locality theorem does and does not identify}

The activities may depend on the whole background-independent
finite-volume baseline $C_{0,\Lambda}$. Their variable background
dependence reads only $A|_H$, but volume consistency is not
claimed. A baseline chosen separately at each arbitrary
background would localize only perturbations from that reference.
The product background domain and the small local precision
norm are substantive hypotheses, not consequences of a weak
global gauge-field norm or of covariance decay alone.

A deterministic path choice may preserve translations while
breaking discrete rotations; arbitrary decompositions do not
make each activity gauge invariant. An application needing
these symmetries must supply an equivariant construction or
an admissible averaging argument.
The periodic extension later in this notebook verifies uniform
volume counting and the required path/hull construction with the
same length and cardinality bounds, including short tori.
No actual regulator coordinates, source coefficients, moving
gauge chart or admitted history family has been identified by
the abstract hypotheses alone.
